%% ============================================================
%% 模糊综合评价段落模板
%% 变量:
%%   n_criteria         — 指标数量
%%   criteria_names     — 指标名称列表
%%   grade_names        — 评语等级列表 (如 ["优","良","中","差"])
%%   n_grades           — 等级数量
%%   weights            — 权重向量
%%   weights_source     — 权重来源说明
%%   R_matrix           — 模糊评价矩阵 (二维列表)
%%   B_vector           — 综合评价向量
%%   final_grade        — 最终评价等级
%%   membership_fig_path — 隶属度分布图路径 (可选)
%% ============================================================

\subsection{模糊综合评价模型}

\subsubsection{模型原理}

模糊综合评价是基于模糊数学隶属度理论，将定性评价转化为定量分析的方法，特别适合处理"模糊性"和"主观性"较强的评价问题。

\textbf{步骤一：确定因素集与评语集。}
\begin{itemize}
    \item 因素集：$U = \{<< criteria_names | join(",\\ ") >>\}$，共 $n = << n_criteria >>$ 个因素。
    \item 评语集：$V = \{<< grade_names | join(",\\ ") >>\}$，共 $<< n_grades >>$ 个等级。
\end{itemize}

\textbf{步骤二：构造模糊评价矩阵 $\mathbf{R}$。} 通过专家评分或隶属函数，确定各因素对各评语等级的隶属度，得到 $\mathbf{R} = (r_{ij})_{n \times << n_grades >>}$。

\textbf{步骤三：确定权重向量。} 本文采用<< weights_source >>确定权重向量 $\mathbf{W} = (<< weights | map("%.4f" | format) | join(",\\ ") >>)$。

\textbf{步骤四：模糊合成运算。}
\begin{equation}
    \mathbf{B} = \mathbf{W} \circ \mathbf{R} = (b_1, b_2, \ldots, b_{<< n_grades >>})
    \label{eq:fuzzy_composition}
\end{equation}
此处 "$\circ$" 表示模糊合成算子（本文采用加权平均型 $M(\cdot, +)$ 算子）。

\textbf{步骤五：根据最大隶属度原则确定评价等级。}

\subsubsection{模糊评价矩阵}

\begin{table}[H]
    \centering
    \caption{模糊评价矩阵 $\mathbf{R}$}
    \label{tab:fuzzy_R}
    \begin{tabular}{l<% for g in grade_names %>c<% endfor %>}
        \toprule
        \textbf{指标} & <% for g in grade_names %><% if not loop.first %> & <% endif %>\textbf{<< g >>}<% endfor %> \\
        \midrule
<% for i in range(n_criteria) %>
        << criteria_names[i] >> & <% for j in range(n_grades) %><% if not loop.first %> & <% endif %><< "%.4f" | format(R_matrix[i][j]) >><% endfor %> \\
<% endfor %>
        \bottomrule
    \end{tabular}
\end{table}

\subsubsection{综合评价结果}

模糊合成运算得到综合评价向量：
\begin{equation}
    \mathbf{B} = (<< B_vector | map("%.4f" | format) | join(",\\ ") >>)
\end{equation}

根据最大隶属度原则，$\max(b_k) = << "%.4f" | format(max(B_vector)) >>$，对应评语等级为"\textbf{<< final_grade >>}"。

<% if membership_fig_path %>
\begin{figure}[H]
    \centering
    \includegraphics[width=0.65\textwidth]{<< membership_fig_path >>}
    \caption{综合评价向量隶属度分布}
    \label{fig:fuzzy_membership}
\end{figure}
<% endif %>